\documentclass[a4paper,12pt,twoside]{memoir}

% Castellano
\usepackage[spanish,es-tabla]{babel}
\selectlanguage{spanish}
\usepackage{graphicx}
\usepackage{float}
\usepackage[utf8]{inputenc}
\usepackage[T1]{fontenc}
\usepackage{lmodern} % scalable font
\usepackage{microtype}
\usepackage{placeins}

\usepackage{url}

\RequirePackage{booktabs}
\RequirePackage[table]{xcolor}
\RequirePackage{xtab}
\RequirePackage{multirow}

% Links
\PassOptionsToPackage{hyphens}{url}\usepackage[colorlinks]{hyperref}
\hypersetup{
	allcolors = {red}
}

% Ecuaciones
\usepackage{amsmath}

% Rutas de fichero / paquete
\newcommand{\ruta}[1]{{\sffamily #1}}

% Párrafos
\nonzeroparskip

% Huérfanas y viudas
\widowpenalty100000
\clubpenalty100000

% Evitar solapes en el header
\nouppercaseheads

% Imagenes
\usepackage{graphicx}
\newcommand{\imagen}[2]{
	\begin{figure}[!h]
		\centering
		\includegraphics[width=0.9\textwidth]{#1}
		\caption{#2}\label{fig:#1}
	\end{figure}
	\FloatBarrier
}

\newcommand{\imagenflotante}[2]{
	\begin{figure}%[!h]
		\centering
		\includegraphics[width=0.9\textwidth]{#1}
		\caption{#2}\label{fig:#1}
	\end{figure}
}



% El comando \figura nos permite insertar figuras comodamente, y utilizando
% siempre el mismo formato. Los parametros son:
% 1 -> Porcentaje del ancho de página que ocupará la figura (de 0 a 1)
% 2 --> Fichero de la imagen
% 3 --> Texto a pie de imagen
% 4 --> Etiqueta (label) para referencias
% 5 --> Opciones que queramos pasarle al \includegraphics
% 6 --> Opciones de posicionamiento a pasarle a \begin{figure}
\newcommand{\figuraConPosicion}[6]{%
  \setlength{\anchoFloat}{#1\textwidth}%
  \addtolength{\anchoFloat}{-4\fboxsep}%
  \setlength{\anchoFigura}{\anchoFloat}%
  \begin{figure}[#6]
    \begin{center}%
      \Ovalbox{%
        \begin{minipage}{\anchoFloat}%
          \begin{center}%
            \includegraphics[width=\anchoFigura,#5]{#2}%
            \caption{#3}%
            \label{#4}%
          \end{center}%
        \end{minipage}
      }%
    \end{center}%
  \end{figure}%
}

%
% Comando para incluir imágenes en formato apaisado (sin marco).
\newcommand{\figuraApaisadaSinMarco}[5]{%
  \begin{figure}%
    \begin{center}%
    \includegraphics[angle=90,height=#1\textheight,#5]{#2}%
    \caption{#3}%
    \label{#4}%
    \end{center}%
  \end{figure}%
}
% Para las tablas
\newcommand{\otoprule}{\midrule [\heavyrulewidth]}
%
% Nuevo comando para tablas pequeñas (menos de una página).
\newcommand{\tablaSmall}[5]{%
 \begin{table}
  \begin{center}
   \rowcolors {2}{gray!35}{}
   \begin{tabular}{#2}
    \toprule
    #4
    \otoprule
    #5
    \bottomrule
   \end{tabular}
   \caption{#1}
   \label{tabla:#3}
  \end{center}
 \end{table}
}

%
%Para el float H de tablaSmallSinColores
\usepackage{float}

%
% Nuevo comando para tablas pequeñas (menos de una página).
\newcommand{\tablaSmallSinColores}[5]{%
 \begin{table}[H]
  \begin{center}
   \begin{tabular}{#2}
    \toprule
    #4
    \otoprule
    #5
    \bottomrule
   \end{tabular}
   \caption{#1}
   \label{tabla:#3}
  \end{center}
 \end{table}
}

\newcommand{\tablaApaisadaSmall}[5]{%
\begin{landscape}
  \begin{table}
   \begin{center}
    \rowcolors {2}{gray!35}{}
    \begin{tabular}{#2}
     \toprule
     #4
     \otoprule
     #5
     \bottomrule
    \end{tabular}
    \caption{#1}
    \label{tabla:#3}
   \end{center}
  \end{table}
\end{landscape}
}

%
% Nuevo comando para tablas grandes con cabecera y filas alternas coloreadas en gris.
\newcommand{\tabla}[6]{%
  \begin{center}
    \tablefirsthead{
      \toprule
      #5
      \otoprule
    }
    \tablehead{
      \multicolumn{#3}{l}{\small\sl continúa desde la página anterior}\\
      \toprule
      #5
      \otoprule
    }
    \tabletail{
      \hline
      \multicolumn{#3}{r}{\small\sl continúa en la página siguiente}\\
    }
    \tablelasttail{
      \hline
    }
    \bottomcaption{#1}
    \rowcolors {2}{gray!35}{}
    \begin{xtabular}{#2}
      #6
      \bottomrule
    \end{xtabular}
    \label{tabla:#4}
  \end{center}
}

%
% Nuevo comando para tablas grandes con cabecera.
\newcommand{\tablaSinColores}[6]{%
  \begin{center}
    \tablefirsthead{
      \toprule
      #5
      \otoprule
    }
    \tablehead{
      \multicolumn{#3}{l}{\small\sl continúa desde la página anterior}\\
      \toprule
      #5
      \otoprule
    }
    \tabletail{
      \hline
      \multicolumn{#3}{r}{\small\sl continúa en la página siguiente}\\
    }
    \tablelasttail{
      \hline
    }
    \bottomcaption{#1}
    \begin{xtabular}{#2}
      #6
      \bottomrule
    \end{xtabular}
    \label{tabla:#4}
  \end{center}
}

%
% Nuevo comando para tablas grandes sin cabecera.
\newcommand{\tablaSinCabecera}[5]{%
  \begin{center}
    \tablefirsthead{
      \toprule
    }
    \tablehead{
      \multicolumn{#3}{l}{\small\sl continúa desde la página anterior}\\
      \hline
    }
    \tabletail{
      \hline
      \multicolumn{#3}{r}{\small\sl continúa en la página siguiente}\\
    }
    \tablelasttail{
      \hline
    }
    \bottomcaption{#1}
  \begin{xtabular}{#2}
    #5
   \bottomrule
  \end{xtabular}
  \label{tabla:#4}
  \end{center}
}



\definecolor{cgoLight}{HTML}{EEEEEE}
\definecolor{cgoExtralight}{HTML}{FFFFFF}

%
% Nuevo comando para tablas grandes sin cabecera.
\newcommand{\tablaSinCabeceraConBandas}[5]{%
  \begin{center}
    \tablefirsthead{
      \toprule
    }
    \tablehead{
      \multicolumn{#3}{l}{\small\sl continúa desde la página anterior}\\
      \hline
    }
    \tabletail{
      \hline
      \multicolumn{#3}{r}{\small\sl continúa en la página siguiente}\\
    }
    \tablelasttail{
      \hline
    }
    \bottomcaption{#1}
    \rowcolors[]{1}{cgoExtralight}{cgoLight}

  \begin{xtabular}{#2}
    #5
   \bottomrule
  \end{xtabular}
  \label{tabla:#4}
  \end{center}
}




\graphicspath{ {./img/} }

% Capítulos
\chapterstyle{bianchi}
\newcommand{\capitulo}[2]{
	\setcounter{chapter}{#1}
	\setcounter{section}{0}
	\setcounter{figure}{0}
	\setcounter{table}{0}
	\chapter*{#2}
	\addcontentsline{toc}{chapter}{#2}
	\markboth{#2}{#2}
}

% Apéndices
\renewcommand{\appendixname}{Apéndice}
\renewcommand*\cftappendixname{\appendixname}

\newcommand{\apendice}[1]{
	%\renewcommand{\thechapter}{A}
	\chapter{#1}
}

\renewcommand*\cftappendixname{\appendixname\ }

% Formato de portada
\makeatletter
\usepackage{xcolor}
\newcommand{\tutor}[1]{\def\@tutor{#1}}
\newcommand{\course}[1]{\def\@course{#1}}
\definecolor{cpardoBox}{HTML}{E6E6FF}
\def\maketitle{
  \null
  \thispagestyle{empty}
  % Cabecera ----------------
\noindent\includegraphics[width=\textwidth]{cabecera}\vspace{1cm}%
  \vfill
  % Título proyecto y escudo informática ----------------
  \colorbox{cpardoBox}{%
    \begin{minipage}{.8\textwidth}
      \vspace{.5cm}\Large
      \begin{center}
      \textbf{TFG del Grado en Ingeniería Informática}\vspace{.6cm}\\
      \textbf{\LARGE\@title{}}
      \end{center}
      \vspace{.2cm}
    \end{minipage}

  }%
  \hfill\begin{minipage}{.20\textwidth}
    \includegraphics[width=\textwidth]{escudoInfor}
  \end{minipage}
  \vfill
  % Datos de alumno, curso y tutores ------------------
  \begin{center}%
  {%
    \noindent\LARGE
    Presentado por \@author{}\\ 
    en Universidad de Burgos --- \@date{}\\
    Tutor: \@tutor{}\\
  }%
  \end{center}%
  \null
  \cleardoublepage
  }
\makeatother


% Datos de portada
\title{eLearningQA \\Documentación Técnica}
\author{Alberto Díaz Álvearez}
\tutor{Raúl Marticorena Sánchez y Carlos López Nozal}
\date{\today}

\begin{document}

\maketitle



\cleardoublepage



%%%%%%%%%%%%%%%%%%%%%%%%%%%%%%%%%%%%%%%%%%%%%%%%%%%%%%%%%%%%%%%%%%%%%%%%%%%%%%%%%%%%%%%%



\frontmatter


\clearpage

% Indices
\tableofcontents

\clearpage

\listoffigures

\clearpage

\listoftables

\clearpage

\mainmatter

\appendix

\apendice{Plan de Proyecto Software}

\section{Introducción}
Con el fin de abordar el plan del proyecto, desarrollaremos una guía detallada que incluirá los puntos clave y los avances logrados en cada sprint. Además, se registrarán las herramientas seleccionadas, las horas invertidas y otros aspectos relevantes. En cada sprint, se especificarán las fechas de inicio y finalización, los objetivos propuestos y los objetivos cumplidos, junto con los problemas encontrados durante el proceso, lo que nos permitirá mostrar la evolución que hemos logrado hasta el momento.

Además, dentro de la guía del plan de proyecto, abordaremos los requisitos detallados del mismo. Esto incluirá un análisis exhaustivo de las funcionalidades y características que la aplicación debe cumplir para satisfacer las necesidades del cliente y los usuarios finales.

También se contemplará el diseño del proyecto, donde se describirán los aspectos técnicos y estéticos que guiarán el desarrollo.

Además, se proporcionará un manual específico destinado al programador, donde se detallarán las pautas y las mejores prácticas para la codificación de la aplicación. Este manual incluirá información sobre los lenguajes de programación utilizados, las convenciones de nomenclatura, las directrices de estilo de codificación y cualquier otro detalle técnico que facilite el desarrollo del proyecto de manera eficiente y coherente.

Asimismo, se creará un manual orientado al usuario final. Este manual contendrá instrucciones claras y concisas sobre cómo utilizar la aplicación o proyecto una vez finalizado. Se incluirán descripciones detalladas de las funcionalidades, los flujos de trabajo recomendados, las configuraciones disponibles y cualquier otra información relevante para que los usuarios finales puedan aprovechar al máximo la aplicación o proyecto.
\section{Planificación temporal}
En esta sección se estimará el tiempo de trabajo en el TFG tanto en la elaboración de la memoria y anexo como en la del propio programa, detallando los acontecimientos a lo largo del proceso.

La planificación del proyecto se ha llevado a cabo en sprints de una duración de dos semanas. Con una reunión al final de cada sprint para comentar los objetivos del último sprint y los nuevos objetivos, en varias ocasiones fue necesario reducir la duración de dichos sprints a una semana. Para el control de versiones tanto de las memorias como del programa se ha utilizado github con un seguimiento a través de commits dejando un histórico bastante detallado. Además de realizar una release cada vez que se tuviese un prototipo funcional de la aplicación.

A continuación mostraré un resumen de dichos sprints:

\subsection{Sprint 0 (23/02/2023 - 08/03/2023)}
	En esta primera reunión se introdujo el proyecto anterior, posibles herramientas con las que trabajar, recursos necesarios y disponibilidad tanto del alumno como de los tutores. Como objetivos se marcaron: inspeccionar el código ya existente, poner en funcionamiento el anterior proyecto, realizar pequeños cambios en el código, inspeccionar el WebService de moodle, instalación y creación de moodle en localhost para empezar a trabajar en local y crear el esqueleto de los documentos en LaTex.
    También añado a la reunión de este sprint que el profesor Carlos López se encontraba de baja por lo que esta y las posteriores reuniones hasta su incorporación las haría con Raúl Marticorena.

\subsection{Sprint 1 (08/03/2023 - 22/03/2023)}
	Durante el segundo sprint tuve varios problemas tanto para entender cómo funcionaba el código como para hacerlo funcionar ya que Moodle había cambiado a la versión 4.0 con cambios en sus respuesta en las llamadas de los web services. Un ejemplo es el paso de campos como "isVisible" de un entero a un booleano. Dados estos problemas y mi falta de tiempo el profesor Raúl Marticorena y yo indicamos como objetivos seguir con la refactorización, enviar el repositorio de github, probar la aplicación contra mountorange e investigar la integración con Heroku aunque este último punto era el menos prioritario para nosotros.

\subsection{Sprint 2 (22/03/2023 - 31/03/2023)}
	Este sprint se decidió hacer de una semana porque a las dos semanas era semana santa y vimos más prioritario realizar la reunión antes y no aplazar la reunión una semana más dejando el sprinte de tres o cuatro semanas. Durante esta semana conseguí que la aplicación se adaptase completamente a la nueva versión de Moodle. Para el siguiente sprint seguiamos con los mismos objetivos ya que no había avanzado mucho.

\subsection{Sprint 3 (31/03/2023 - 17/04/2023)}
	En la revisión de este sprint mostré la incorporación de los modelos necesarios para poder trabajar con cuestionarios en la aplicación y descubrí que Heroku se había vuelto de pago incluso para las aplicaciones de prueba por lo que dejamos el tema de buscra un servidor online de lado. Se mantienen objetivos para el siguiente sprint y se deja la próxima reunión para dentro de una semana.

\subsection{Sprint 4 (17/04/2023 - 27/04/2023)}
	Para esta reunión se había inorporado el un campo en los informes que mostraba si el curso seleccionado tenía cuestionarios o no. También se añadió un registro que mostraba el porcentaje de alumnos que realizaban los cuestionarios del curso. Esta es la primera vez del trabajo que realizo una release ya que teníamos algo avanzada la aplicación. Para el siguiente sprint decidimos mantener la reunión semanal y pusimos como objetivo realizar un resumen del cuestionario al hacer click en dicho cuestionario junto con algún gráfico.

\subsection{Sprint 5 (27/04/2023 - 03/05/2023)}
	Durante este sprint avancé en la lógica de obtención de datos de dicho resumen pero al llegar la reunión no tenía ninguna implementación definitiva por lo que no se mostraban avances en la interfaz. Para la semana siguiente debería tener implementado ya dicho resumen.

\subsection{Sprint 6 (03/05/2023 - 11/05/2023)}
	Para este sprint ya aparecía un resumen en el cuestionario mostrando datos básicos a falta del gráfico. El profesor Raúl MArticorena vió importante meter alguna estadística como el coeficiente de curtosis (con este valor podríamos saber si hay valores atípicos en las notas) y la asimetría. Para la siguiente semana dejamos como objetivo implementar estos datos y el gráfico.

\subsection{Sprint 7 (11/05/2023 - 18/05/2023)}
	Durante esta semana terminé primero los datos del resumen restantes aunque Raúl Marticorena dijo que se estaban calculando mal por lo que para la siguiente semana tendré que corregir la forma en la que los calcula y respecto al gráfico se implementó la lógica pero el envío de la información de Java a Javascript no tuvo éxito por lo que para la siguiente reunión dejamos como objetivo terminar dicho gráfico, empezar a desarrollar la memoria y comenzar con la parte de foros.

\subsection{Sprint 8 (18/05/2023 - 25/05/2023)}
	En esta reunión se incorporó el profesor Carlos López por lo que los primeros minutos fueron para ver el progreso en el proyecto, después aporto una serie de mejoras y fallos que no habíamos contemplado. Los cambios que se implementaron durante este sprint fueron principalmente arreglar los métodos que realizaban cálculos sobre las estadísticas de los cuestionarios, se consiguió implementar correctamente el gráfico (a falta de normalizar el valor de cada pregunta ya que se llegó a la conclusión de que habría preguntas de diferente puntuación) y se generó una release con el tema de los cuestionarios terminado. La parte de los foros se aplazó al siguiente sprint por falta de tiempo ya que quedaban apenas dos semanas para los exámenes finales.

\subsection{Sprint 9 (25/05/2023 - 02/06/2023)}
	Durante este sprint se pulieron algunos fallos detectados por los profesores como la normalización de las notas, el número de intentos y agregar el formato de notas en español (se usan comas en vez de puntos). Luego se agregaron los porcentajes de participación en los foros, se redujo el número de llamadas a la API de Moodle. También se avanzo en la redacción de los anexos. En esta reunión se decidió que si queríamos llegar a tiempo tendría que formalizarse un borrador de la memoria, de los anexos, generar una release totalmente funcional con todos los cambios e implementar los actions de github como es el de ejecutar SonarCloud cada vez que se haga un push a la rama develop.

 \subsection{Sprint 10 (02/06/2023 - 08/06/2023)}
	Esta reunión era clave para ver si se hacía entrega del trabajo en primera convocatoria o en segunda. Al ver que se había desarrollado una memoria y unos anexos que estaban muy verdes y sin estar revisados además de comprobar que la aplicación daba algunos fallos se decidió aplazar a segunda convocatoria. También se decidió poner un poco de organización y aunque tarde pero se empezó a hacer uso de las issues y milestones. De ahora en adelante el objetivo era corregir errores, redactar una memoria de calidad e intentar incluir unas últimas mejoras que harían de este un trabajo más completo.

 \subsection{Sprint 11 (08/06/2023 - 26/06/2023)}
	Durante este sprint que fue un poco más largo se corrigieron errores en los cálculos de las estadísticas tanto en los cuestionarios como en los foros, se introdujo una página de carga previa a la muestra del informe ya que el tiempo de carga cada vez era mayor. En la reunión se habló de aonsejar trabajar la asincronía de la aplicación en futuras versiones. Con la aplicación funcionando correctamente comenzó el enfoque en la memoria para el siguiente sprint y en la incorporación de análisis de sentimientos en los foros.

\section{Estudio de viabilidad}

\subsection{Viabilidad económica}
En este apartado se estiman los costes que tendría la realización del proyecto en un caso real. Como los cambios que se han hecho han sido implementaciones al proyecto ya creado se mantendrá la información del compañero de la primera versión\cite{previotfganexos} y se resaltarán en las tablas los cambios con otro color para comprender qué se ha modificado.

La primera versión del proyecto supuso 8 meses de desarrollo con unas 800 horas de trabajo, esto equivaldría a 5 meses. La segunda versión que detalla el tema de los cuestionarios y foros supone en torno a unas 320 horas que equivale a 2 meses de trabajo.
\begin{table}[H]
	\centering
	\caption{Costes de personal}
	\resizebox{0.5\textwidth}{!}{%
		\begin{tabular}{lc}
			\hline
			\textbf{Concepto}        & \textbf{Coste} \\ \hline
			Sueldo mensual neto      & 1.000,00€      \\
			Retención IRPF (19\%)    & 254,52€        \\
			Seguridad Social (6,35\%) & 85,06€         \\
			Sueldo mensual bruto     & 1.339,58€      \\ \hline
			\textbf{\textcolor{blue}{Total 7 meses}}   & \textcolor{blue}{9.377,06€}      \\ \hline
		\end{tabular}%
	}
\end{table}
Las cuotas a la seguridad social se componen de un 4,70\% de contingencias comunes, un 1,55\% por desempleo de tipo general, y un 0,10\% de formación profesional.

Se supone que el ordenador portátil se amortiza en cinco años y ha sido usado cinco meses.
\begin{table}[H]
	\centering
	\caption{Costes de hardware}
	\resizebox{0.75\textwidth}{!}{%
		\begin{tabular}{lcc}
			\hline
			\textbf{Concepto}        & \textbf{Coste}& \textbf{Coste amortizado} \\ \hline
			Ordenador portátil     & 550€ & 45,83€    \\ \hline
			\textbf{Total}   & 550€ & 45,83€    \\ \hline
		\end{tabular}%
	}
\end{table}

\begin{table}[H]
	\centering
	\caption{Costes varios}
	\resizebox{0.5\textwidth}{!}{%
		\begin{tabular}{lc}
			\hline
			\textbf{Concepto}        & \textbf{Coste} \\ \hline
			Internet      & 150,00€      \\
			Electricidad    & 175,00€        \\ \hline
			\textbf{Total}   & 325,00€      \\ \hline
		\end{tabular}%
	}
\end{table}

A partir de estos costes obtenemos el coste total del proyecto:
\begin{table}[H]
	\centering
	\caption{Costes totales}
	\resizebox{0.5\textwidth}{!}{%
		\begin{tabular}{lc}
			\hline
			\textbf{Concepto}        & \textbf{Coste} \\ \hline
			Personal      & \textcolor{blue}{9.377,06€}      \\
			Hardware    & 45,83€           \\
			Varios     & 325,00€      \\ \hline
			\textbf{\textcolor{blue}{Total 7 meses}}    & \textcolor{blue}{9.747,89€}      \\ \hline
		\end{tabular}%
	}
\end{table}

Para rentabilizar la aplicación web se podrían tener en cuenta varias posibles soluciones como ofrecer un sistema freenium dejando parte del contenido de forma gratuita mientras que otras otros apartados serían premium ocasionando un coste como es el caso de Spotify. Otra posible solución es la que comenta mi compañero Roberto en la versión previa mencionando una suscripción menusal los que implicaría unos ingresos continuos\cite{previotfganexos}.

\subsection{Viabilidad legal}
\subsubsection{Licencias de software}
Respecto a la viabilidad legal en esta segunda versión solo se ha incorporado GitHub Copilot el cual no ofrece ningún problema legal ya que sus suscripción mensual permite el libre manejo del código generado. A continuación se mostrarán los frameworks y librerías con sus respectivas licencias.

\begin{table}[H]
	\caption{Licencias del software utilizado}
	\resizebox{\textwidth}{!}{%
		\begin{tabular}{|c|c|c|}
			\hline
			\textbf{Software}   & \textbf{Descripción}                                                                                                         & \textbf{Licencia} \\ \hline
			Spring Framework    & Framework para aplicaciones web                                                                                              & Apache 2.0        \\ \hline
			Tomcat Embed Jasper & \begin{tabular}[c]{@{}c@{}}Implementación de Tomcat que incluye\\ Jasper, el parser de JSP de Tomcat\end{tabular}            & Apache 2.0        \\ \hline
			JUnit               & Framework para tests unitarios en Java                                                                                       & EPL               \\ \hline
			Apache Commons IO   & \begin{tabular}[c]{@{}c@{}}Librería de utilidades varias (usado en \\ traducción de imágenes a arrays de bytes)\end{tabular} & Apache 2.0        \\ \hline
			Apache Log4j        & Librería para registro de logs                                                                                               & Apache 2.0        \\ \hline
			Bootstrap           & \begin{tabular}[c]{@{}c@{}}Librerías CSS y JavaScript para \\ páginas web\end{tabular}                                       & MIT               \\ \hline
			Plotly.js           & \begin{tabular}[c]{@{}c@{}}Librería JavaScript de \\ generación de gráficos\end{tabular}                                     & MIT               \\ \hline
			GitHub Copilot           & \begin{tabular}[c]{@{}c@{}}GitHub Copilot es una\\ herramienta de inteligencia artificial\end{tabular}                                     & GitHub, Inc.               \\ \hline
		\end{tabular}%
	}
\end{table}
La licencia pública de Eclipse (EPL) es compatible con estas licencias sin contar GitHub Copilot, a continuación se mencionan sus posibilidades y obligaciones.
\begin{itemize}
	\item\textbf{Permite:} uso, reproducción, distribución, modificación, uso comercial y uso de patentes.
	\item\textbf{Obliga a:} revelar la fuente y el autor, mantener la misma licencia al redistribuir el software, distribuir el software libre de regalías.
	\item\textbf{No permite:} responsabilizar al autor o contribuidores por posibles daños, utilizar marcas propiedad del autor para promoción o publicidad.
\end{itemize}

GitHub Copilot es una herramienta de programación desarrollada por GitHub y OpenAI. Utiliza la tecnología de inteligencia artificial de GPT-3 para ofrecer sugerencias de código en tiempo real mientras escribes. Copilot analiza el contexto del código que estás escribiendo y genera automáticamente fragmentos de código relevantes y útiles para completar tus líneas de código.
Tiene 2 tarifas:
\begin{itemize}
	\item 10 dólares al mes para uso personal.
	\item 19 dólares al mes por persona para uso empresarial.
\end{itemize}
\apendice{Especificación de Requisitos}
\label{apendice:B}
\section{Introducción}

\section{Objetivos generales}
	
\apendice{Especificación de diseño}

\section{Introducción}
En este apéndice se proporciona una descripción detallada de cómo se tiene previsto cumplir con el conjunto de objetivos y requisitos establecidos en el apéndice \ref{apendice:B}. Se incluye información sobre los datos utilizados por la aplicación, su arquitectura y su diseño de procedimientos.
\section{Diseño de datos}
En la siguiente tabla se recogen las entidades asociadas a las consultas incorporadas en la nueva versión.
\begin{table}[H]
	\resizebox{\textwidth}{!}{%
		\begin{tabular}{|l|l|}
			\hline
			\textbf{Consulta}                                                                                                                                        & \textbf{Entidades}        \\ \hline
			El curso tiene al menos un cuestionario                                                                                                         & Curso, Cuestionario            \\ \hline
			El curso tiene al menos un foro                                                                                                                 & Curso, Foro            \\ \hline
			Un mínimo de alumnos participa en los cuestionarios                                                                                             & Curso, Cuestionario            \\ \hline
			Un mínimo de alumnos participa en los foros                                                                                                     & Curso, Foro            \\ \hline
		\end{tabular}%
	}
\end{table}
Seguidamente se mostrará una imagen del modelo de datos necesario para comprender la estructura con la que trabaja Moodle y los datos que se deberán recuperar.
\imagen{ModeloDatos.png}{Diagrama de clases de los modelos utilizados en la segunda versión}
\section{Diseño procedimental}
A continuación, se mostrará un diagrama del flujo que implementa la aplicación dividiéndose en las pantallas de login, lista de cursos e informe. La dirección de las flechas indica que al hacer login se accede a la lista de cursos y al hacer logout se vuelve a la pantalla de login, dentro de la pantalla con la lista de cursos se indica que se puede acceder también a la pantalla de los informes. Estas interacciones se detallarán más adelante en el trabajo.
\imagen{BocetoInterfaz.png}{Boceto de la interfaz}
\subsection{Patrón fachada}
El patrón de diseño Fachada es un patrón estructural que proporciona una interfaz simplificada para interactuar con un sistema complejo o conjunto de clases. La idea principal detrás del patrón Fachada es proporcionar una interfaz de nivel superior que simplifique el uso del sistema subyacente. En lugar de que los clientes interactúen directamente con múltiples componentes o clases internas, pueden hacerlo a través de la fachada, que se encarga de manejar las interacciones internas y ofrecer una interfaz más simple y específica.
\imagen{DiagramaFachada.png}{Diagrama de interacción con la fachada}
A continuación se mostrará la definición de Roberto Arasti de la primera versión la cuál no varía y solamente cabría destacar que en este trabajo el patrón de diseño Fachada se ha utilizado para solicitar a la clase WebServiceClient las estadísticas de cuestionarios y foros para luego mostrarlos en el jsp generado.
Seguidamente, se explicará como acceder a la información a través del patrón de diseño fachada:
cuando un usuario intenta acceder a una página específica, como por ejemplo el informe de la aplicación, la clase SpringController se encarga de manejar la solicitud y proporcionar la respuesta correspondiente. En este caso, la respuesta es un archivo JSP, el cual debe ser compilado y ejecutado para generar una página HTML como resultado final. Al ejecutar el archivo JSP en cuestión, se desencadenan las llamadas necesarias a la fachada para realizar las comprobaciones predefinidas. Estas comprobaciones se llevan a cabo a través de la clase WebServiceClient, que accede al servidor de Moodle. Después de recibir los resultados de las comprobaciones, el archivo JSP utiliza la fachada como una clase de utilidad para procesar algunos de los resultados. Una vez finalizada la ejecución del archivo JSP, la clase SpringController devuelve la página HTML generada al usuario \cite{previotfganexos}.
\apendice{Documentación técnica de programación}

\section{Introducción}
En este apéndice se proporciona una guía detallada destinada a futuros programadores que utilicen el proyecto, con el objetivo de brindarles orientación y ayuda. Se incluye información sobre la estructura de directorios del proyecto, donde se explican en detalle los contenidos de cada uno. Además, se presenta un manual del programador que aborda aspectos relevantes que podrían resultar útiles. También se detallan los pasos necesarios para compilar, instalar y ejecutar el proyecto en un entorno de desarrollo integrado. Por último, se describen las pruebas del sistema que se han llevado a cabo, ofreciendo así una visión completa de las pruebas realizadas.
\section{Estructura de directorios}
La estructura del repositorio se detalla a continuación:
\begin{itemize}
	\item \textbf{/}: contiene en archivo README que hace una breve explicación del funcionamiento de la aplicación y las carpetas contenedoras de la aplicación y los ficheros de la memoria.
	\item \textbf{/.github/workflows}: dirección de los actions que serán ejecutados en GitHub, en este caso será un archivo maven.yml. Este archivo indica que cada vez que se haga un push a la rama develop se realizará un análisis de los últimos cambios en SonarCloud.
	\item \textbf{/Memo/plantillaLatex-master}: dirección de la memoria, anexos y archivos en formato \LaTeX para la configuración de los distintos PDFs.
	\item \textbf{/Memo/plantillaLatex-master/img}: carpeta dónde se alojan las imágenes que utilizarán la memoria y los anexos.
	\item \textbf{/Memo/plantillaLatex-master/tex}: ficheros en formato \LaTeX que sirven para generar los PDFs de la memoria y anexos.
	\item \textbf{/Project/prototipo}: carpeta que contiene la aplicación completa.
	\item \textbf{/Project/prototipo/configurations}: Carpeta con los archivos que asignan las propiedades del proyecto.
	\item \textbf{/Project/prototipo/src}: código fuente de la aplicación Java.
	\item \textbf{/Project/prototipo/src/main}: este directorio abarca el código fuente de la aplicación, los casos de prueba y los recursos empleados en su funcionamiento. 
	\item \textbf{/Project/prototipo/src/main/java/es/ubu/lsi}: Dirección de las clases que permiten el correcto funcionamiento de la aplicación web.
	\item \textbf{/Project/prototipo/src/main/java/es/ubu/lsi/model}: dentro de este directorio se encuentran las clases Java Bean necesarias para la deserialización de los datos obtenidos a través de servicios web. Estas clases permiten transformar los datos extraídos en objetos que pueden ser fácilmente manipulados y utilizados en el desarrollo de la aplicación.
	\item \textbf{/Project/prototipo/src/main/resources}: contiene el archivo de propiedades de la aplicación web, el cual contiene configuraciones y ajustes específicos para su correcto funcionamiento.
	\item \textbf{/Project/prototipo/src/main/resources/images}: ubicación de las imágenes desplegadas en la web.
	\item \textbf{/Project/prototipo/src/main/resources/json}: información de los cursos que se utilizarán en los casos de pruebas.
	\item \textbf{/Project/prototipo/src/main/resources/json/informe}: ubicación donde se encuentran los textos estáticos utilizados en la tabla del informe de fases, los cuales están almacenados en formato JSON.
	\item \textbf{/Project/prototipo/src/main/resources/static/js}: librerías JavaScript que utiliza la aplicación.
	\item \textbf{/Project/prototipo/src/main/webapp/WEB-INF/jsp}: contiene las páginas de la aplicación en formato jsp.
	\item \textbf{/Project/prototipo/src/test/java/es/ubu/lsi}: contiene los casos de pruebas sobre las consultas que realiza la aplicación en los cursos.
	\item \textbf{/Project/prototipo/target}: es el directorio donde Maven deposita los resultados, como las clases compiladas y la aplicación compilada en formato WAR.
	
\end{itemize}
\section{Manual del programador}
\subsection{Entorno}
Requisitos del proyecto:
\begin{itemize}
	\item Java JDK 8
	\item Un IDE como puede ser Eclipse, NetBeans o IntelliJ IDEA
\end{itemize}

Para descargar el JDK 8 empleado en la aplicación es necesaria una cuenta de Oracle y posteriormente obtener el JDK en el siguiente enlace \url{https://www.oracle.com/java/technologies/downloads/#java8}.

En este caso como se ha trabajado con eclipse adjunto el enlace de la descarga \url{https://www.eclipse.org/downloads/}

\subsection{Obtención del código fuente}
Para obtener el código fuente de la aplicación web solo vamos a necesitar clonar el proyecto en nuestro equipo o en el caso de eclipse dentro del propio IDE en la versión importar proyecto Git. El enlace al proyecto es \url{https://github.com/ada1012/eLearningQA}. Si tenemos intención de hacer cambios en la aplicación, antes de clonar el proyecto deberemos hacer un fork para así tener una copia en nuestro perfil de dicho proyecto. En mi caso al manejar git con la extensión de Visual Studio Code puedo manejar los pull, push y commits desde la interfaz gráfica.
\imagen{GitVSCode.png}{Manejo de Git con Visual Studio Code}
\subsection{Propósito de las clases principales}
\begin{itemize}
	\item \textbf{Application:} clase principal de la aplicación desde la que se arranca la aplicación web.
	\item \textbf{SpringController:} indica mediante anotaciones y funciones qué devolver cuando el usuario accede a cada página. En este caso solo devuelve páginas jsp e imágenes, pero se podrían devolver otros tipos de recurso, como archivos de audio o documentos.
	\item \textbf{ELearningQAFacade:} hace de intermediario entre el usuario y la clase WebServiceClient.También contiene funciones para generar tablas y sus contenidos a partir de la información obtenida. 
	\item \textbf{WebServiceClient:} contiene el código de las comprobaciones y realiza las llamadas REST al servidor de Moodle para obtener la información necesaria. También trabajará con los datos obtenidos de Moodle para adaptarlos a las estadísticas de cuestionarios y foros.
	\item \textbf{FacadeConfig:} almacena la configuración elegida por el usuario, junto con el host de Moodle, se almacena como atributo en la fachada para poder pasarlo a las funciones del WebServiceClient y así permitir que este último modifique su comportamiento de acuerdo a la configuración establecida.
	\item \textbf{AlertLog:} la clase AlertLog se encarga de generar y almacenar mensajes de alerta. Estos mensajes se pasan como parámetro a las funciones del WebServiceClient con el fin de que puedan ser almacenados en la lista de mejoras del informe.
	\item \textbf{RegistryIO:} se encarga de almacenar y acceder a la información guardada en archivos CSV que se utilizan para generar los gráficos. Además, proporciona una función que permite generar un gráfico utilizando los datos contenidos en un archivo CSV.
	\item \textbf{AnalysisSnapshot:} representa un conjunto de datos almacenados en cada línea de los archivos CSV. Esta clase proporciona funcionalidades que facilitan el almacenamiento y la manipulación de los datos, así como la generación de gráficos a partir de ellos.
\end{itemize}
\subsection{Integración continua}
Dentro del proyecto se encuentra un archivo denominado \texttt{maven.yml}, ubicado en la ruta \texttt{/.github/workflows}. Este archivo contiene una serie de instrucciones que se ejecutan automáticamente cuando se realiza un push a la rama develop del repositorio. A continuación, proporcionaremos una breve explicación de dichas instrucciones. \imagen{MavenYML.png}{Trabajos de integración continua}
En primer lugar, es notable que la máquina virtual utilizada para la integración continua se basa en la versión más actualizada de Ubuntu. Posteriormente, se procede a la instalación de JDK 11, ya que se ha determinado que es necesario para el proceso de integración continua después de la incorporación de SonarCloud. Además, se realiza el almacenamiento en caché de los paquetes de SonarCloud y Maven en la memoria de la máquina virtual.

Por último, se lleva a cabo la construcción y análisis del proyecto, indicando explícitamente la ubicación del archivo \texttt{pom.xml}. Esta especificación se debe a que el directorio raíz del proyecto difiere del directorio raíz del repositorio, lo cual ha ocasionado algunos inconvenientes que abordaremos posteriormente.
\subsection{Herramienta de calidad de código}
Para analizar la calidad del código del proyecto se ha utilizado SonarCloud como se hizo en el trabajo anterior.

La página principal brinda un resumen conciso del estado actual del proyecto en términos de errores, defectos de código, aspectos de seguridad y la cantidad de código duplicado presente. Además, proporciona información sobre el estado del proyecto con respecto al Quality Gate, que evalúa si se cumplen una serie de condiciones predefinidas.
\imagen{SonarCloud.png}{Página del proyecto en SonarCloud en \url{https://sonarcloud.io/project/overview?id=ada1012_eLearningQA}}
Con esta herramienta se consigue una refactorización inmediata ya que cada vez que se suba un fragmento de código al proyecto esta aplicación resumirá los distintos fallos encontrados, haciendo distinción entre los ya existentes y los nuevos.

\section{Compilación, instalación y ejecución del proyecto}
\subsection{Ejecución en local}
Para ejecutar la aplicación en local solo es necesario ejecutar la clase principal: Application, los IDEs mencionados previamente los ejecutarán sin problema al ejecutar la opción \texttt{Run} y se podrá acceder a dicha aplicación a través de la url: \url{http://localhost:8080/}.
\imagen{EjecucionLocalhost.png}{Ejecución en local}

\section{Pruebas del sistema}
Los test se realizan con JUnit y se han mantenido con respecto a la versión anterior ejecutándose a través del archivo ELearningQAFacadeTest y yendo contra unos archivos estáticos en formato JSON los cuales se encuentran en la ruta "resources/json". Ante las nuevas funcionalidades solo habría que modificar el archivo ELearningQAFacadeTest y asignar unos datos estáticos en el archivo JSON. La opción de obtener los datos de prueba de un servidor de momento es una opción poco viable por el tema del tiempo de carga y el aumento de llamadas en el futuro.
\apendice{Documentación de usuario}

\section{Introducción}
En este apéndice se incluyen los detalles de la aplicación que conciernen al usuario, como los requisitos para utilizar la aplicación y un manual de usuario que también existe en formato web y puede ser accedido desde la página de login de la aplicación.
\section{Requisitos de usuarios}
Para utilizar la aplicación como usuario se necesita:
\begin{itemize}
	\item Una cuenta en un servidor de Moodle con rol de profesor en los cursos a analizar
	\item Conexión a internet para acceder al despliegue de la aplicación
\end{itemize}
\section{Instalación}
Como se ha mencionado en temas anteriores con Java ya instalado el único requisito es ejecutar el archivo WAR que viene en el código fuente de la aplicación.
\section{Manual del usuario}

\subsubsection{Pantallas:}

\textbf{Página de login:} la página de login contiene un formulario de
inicio de sesión, un link ``Acerca de'', un link ``Manual de usuario'' que
muestra la versión web de este manual, y un link de contacto.
\imagen{Login.png}{Página de login}

\textbf{Página principal:~}la página principal contiene una cabecera que
indica el nombre de usuario con el que se ha registrado, un botón para
desconectarse, un enlace para generar un informe global de todos los
cursos, y una tabla con una lista de enlaces para generar los informes
específicos de cada curso con una caja de búsqueda.
\imagen{ListaCursos.png}{Página principal}

\textbf{Página de informe:} la página de informe genera el informe
dependiendo de cómo se haya accedido a esta, ya sea mediante el enlace
``Generar informe global'' o por medio de los enlaces a informes de cursos
específicos.~

Los informes contienen los siguientes elementos:

\begin{itemize}
	\item
	\textbf{Puntuaciones por rol y perspectiva:} el informe contiene una
	matriz de roles y perspectivas que da distintas puntuaciones a cada
	uno de los roles (Diseñador, Facilitador, y Proveedor) en cada una de
	las tres perspectivas definidas (Pedagógica, Tecnológica, y
	Estratégica) según los resultados del análisis.
	\item
	\textbf{Gráfico de evolución:} la aplicación almacena estas puntuaciones cada vez que se carga un informe para generar un gráfico interactivo de la evolución de estos valores.
	\imagen{Evolucion.png}{Página de evolución del rendimiento}
	\item
	\textbf{Informe de fases:} el informe contiene también una tabla que
	muestra los resultados en cada una de las comprobaciones del análisis y
	las agrupa en las distintas fases del diseño instruccional (Diseño,
	Implementación, Realización, y Evaluación) además de dar una
	puntuación general. En el caso de los cuestionarios y foros tendrán cada uno una lista desplegable en el que se podrá ver un desglose de los mismos con el porcentaje de participación de cada uno y el porcenaje de participación general.
	\item
	\textbf{Lista de mejoras:} el informe muestra una lista de
	motivos concretos que producen resultados negativos en el análisis
	para subsanarlos o tenerlos en cuenta de cara al futuro. Haciendo clic en las filas de la tabla se puede filtrar la lista de mejoras.
	\item
	\textbf{Informe de estadísticas:} por último, el informe muestra una sección de estadísticas de los cuestionarios y foros, en el caso de los cuestionarios se mostrará además un gráfico con el índice de dificultad de las preguntas y en el de los foros una lista con los hilos del mismo.
	\imagen{InformeFases.png}{Página de informe de fases}
	\imagen{EstadisticasCuestionario.png}{Estadísticas de un cuestionario}
	\imagen{EstadisticasForo.png}{Estadísticas de un foro}
\end{itemize}

El informe especifico muestra los resultados de las comprobaciones en el informe de fases de forma absoluta, es decir, que solo indica si el criterio se cumple o no, sin embargo, en el informe global, se muestra un resultado u otro ~dependiendo de la cantidad de cursos que satisfacen dicho criterio.


\subsubsection{Acciones del usuario:}

\begin{itemize}
	\item
	\textbf{Login:} para acceder a la aplicación son necesarias las
	credenciales de acceso a una cuenta de la plataforma Moodle a la que
	accede la aplicación (en el caso del prototipo es Mount ~Orange
	School). Debe introducir su usuario y contraseña en los campos
	``Username'' y ``Password''. Si quiere vaciar los campos pulse el botón
	``Borrar''. Para acceder a la página principal, pulse el botón ``Entrar''
	tras haber introducido sus credenciales.
	\item
	\textbf{Desconectar:} desde la página principal, si desea finalizar su
	sesión, pulse el botón ``Desconectar''. Esto invalidará sus credenciales
	y le impedirá acceder a la aplicación hasta que se registre de nuevo
	con unas credenciales válidas.
	\item
	\textbf{Generar informe específico:} la página principal muestra una
	tabla con todos los cursos en los que se encuentra matriculado el
	usuario registrado en formato de enlace. Al clicar un enlace, se
	generará un informe en una pestaña aparte del navegador que mostrará
	los resultados del análisis que ha realizado la aplicación sobre el
	curso correspondiente.
	\item
	\textbf{Generar informe global:} en la página principal hay un enlace
	llamado ``Generar informe global''. Al hacer clic sobre este, se
	generará un informe en una pestaña aparte del navegador que mostrará
	un resumen de los análisis de todos los cursos en los que se encuentra
	matriculado el profesor.
\end{itemize}

\subsubsection{Explicación de las comprobaciones de los informes:}

Las siguientes comprobaciones están relacionadas con los roles, fases, y
perspectivas mencionados anteriormente. Los distintos procesos del
diseño instruccional se encuentran divididos en fases, ~con ciertas
perspectivas en mente, y son responsabilidad directa o indirecta de
ciertos roles. Al estar las comprobaciones ligadas a esos procesos se
muestran agrupadas por fases, y después de la explicación se indican los
roles responsables e involucrados, además de las perspectivas
correspondientes.

\textbf{Diseño:}

\begin{itemize}
	\item
	\textbf{Las opciones de progreso del estudiante están activadas}: se
	comprueba que estén habilitadas las opciones de progreso de los
	estudiantes en el curso.~{Responsable:} Diseñador
	{Involucrados:} Facilitador {Perspectivas:} Pedagógica
	\item
	\textbf{Se proporcionan contenidos en diferentes formatos:} se
	comprueba que haya variedad de formatos en los recursos del curso.
	{Responsable:} Diseñador {Involucrados:} Facilitador y
	Proveedor {Perspectivas:} Pedagógica y Tecnológica
	\item
	\textbf{El curso tiene grupos:} se comprueba que existan grupos
	definidos en el curso. {Responsable:} Diseñador
	{Involucrados:} Facilitador y Proveedor {Perspectivas:}
	Pedagógica
	\item
	\textbf{El curso tiene actividades grupales:} se comprueba que existan
	actividades con entrega grupal habilitada en el curso.
	{Responsable:} Diseñador {Involucrados:} Facilitador y
	Proveedor {Perspectivas:} Pedagógica
	\item
	\textbf{Los estudiantes pueden ver las condiciones necesarias para
		completar una actividad:} se comprueba que esté habilitada la opción
	de mostrar las condiciones para completar una actividad en el curso.
	{Responsable:} Diseñador {Involucrados:} Facilitador y
	Proveedor {Perspectivas:} Pedagógica
	\item
	\textbf{Todas las actividades tienen la misma nota máxima en el
		calificador:} se comprueba que exista una consistencia en las notas
	máximas de los items de calificación (tareas, entregas, cuestionarios)
	del curso. {Responsable:} Diseñador {Involucrados:}
	Facilitador y Proveedor {Perspectivas:} Pedagógica
	\textbf{El curso tiene al menos un cuestionario:} se comprueba que exista un cuestionario dentro de la planificación del curso.
    {Responsable:} Diseñador {Involucrados:} Facilitador y Proveedor {Perspectivas:} Pedagógica
	\textbf{El curso tiene al menos un foro:} se comprueba que exista un foro dentro de la planificación del curso.
    {Responsable:} Diseñador {Involucrados:} Facilitador y Proveedor {Perspectivas:} Pedagógica
\end{itemize}

\textbf{Implementación:}

\begin{itemize}
	\item
	\textbf{Los recursos están actualizados:} se comprueba que los
	recursos del curso tengan una fecha de creación reciente.
	{Responsable:} Diseñador {Involucrados:} Facilitador y
	Proveedor {Perspectivas:} Pedagógica y Tecnológica
	\item
	\textbf{Fechas de apertura y cierre de tareas son correctas:} se
	comprueba que las fechas de apertura y cierre de tareas y
	cuestionarios no se solapen de forma erronea con las fechas de inicio
	y fin del curso. {Responsable:} Facilitador
	{Involucrados:} Diseñador y Proveedor {Perspectivas:}
	Pedagógica y Tecnológica
	\item
	\textbf{Se detallan los criterios de evaluación:} se comprueba que
	exista en al menos una actividad una rúbrica o una guía de calificación
	en el curso.~{Responsable:} Diseñador {Involucrados:}
	Facilitador y Proveedor {Perspectivas:} Pedagógica y Tecnológica
	\item
	\textbf{El calificador no tiene demasiado anidamiento:} se comprueba
	que la estructura de las categorías del calificador no sea demasiado
	enrevesada. {Responsable:} Diseñador {Involucrados:}
	Facilitador y Proveedor {Perspectivas:} Pedagógica y Estratégica
	\item
	\textbf{Todos los alumnos están en algún grupo:} se comprueba que cada
	alumno pertenezca a un grupo. {Responsable:} Proveedor
	{Involucrados:} Diseñador {Perspectivas:} Tecnológica y
	Estratégica
\end{itemize}

\textbf{Realización:}

\begin{itemize}
	\item
	\textbf{El profesor responde en los foros dentro del límite de 48
		horas lectivas desde que se plantea la duda:} se comprueba que no
	hayan preguntas por parte de alumnos que estén sin responder en un
	tiempo razonable. {Responsable:} Facilitador
	{Involucrados:} Diseñador y Proveedor {Perspectivas:}
	Pedagógica y Tecnológica
	\item
	\textbf{Se ofrece retroalimentación de las tareas:} se comprueba que
	el profesor deje comentarios en la mayoría de calificaciones que haga.
	{Responsable:} Facilitador {Involucrados:} Diseñador y
	Proveedor {Perspectivas:} Pedagógica y Tecnológica
	\item
	\textbf{Las tareas están calificadas:} se comprueba que no hayan
	entregas de alumnos que hayan pasado una semana sin calificación.
	{Responsable:} Facilitador {Involucrados:} Diseñador y
	Proveedor {Perspectivas:} Pedagógica y Tecnológica
	\item
	\textbf{El calificador muestra cómo ponderan las diferentes tareas:}
	se comprueba que el calificador muestre los pesos de los items de
	calificación. {Responsable:} Facilitador {Involucrados:}
	Diseñador y Proveedor {Perspectivas:} Pedagógica y Tecnológica
	\item
	\textbf{Al menos un X \% de los alumnos responden a los cuesionarios:}
	se comprueba que hay un mínimo de alumnos que interactúan con los cuestionarios propuestos para el desarrollo del curso. {Responsable:} Proveedor {Involucrados:}
	Diseñador y Facilitador {Perspectivas:} Pedagógica
	\item
	\textbf{Al menos un X \% de los alumnos participa en los foros:}
	se comprueba que hay un mínimo de alumnos que interactúan con los foros propuestos para el desarrollo del curso. {Responsable:} Proveedor {Involucrados:}
	Diseñador y Facilitador {Perspectivas:} Pedagógica
\end{itemize}

\textbf{Evaluación:}

\begin{itemize}
	\item
	\textbf{La mayoría de alumnos responden a los feedbacks:} se comprueba
	que no hayan muchos alumnos que no respondan a los feedbacks.
	{Responsable:} Proveedor {Involucrados:} Diseñador y
	Facilitador {Perspectivas:} Pedagógica, Tecnológica, y
	Estratégica
	\item
	\textbf{Se utilizan encuestas de opinión:} se comprueba que el curso
	contenga encuestas de opinión. {Responsable:} Proveedor
	{Involucrados:} Diseñador y Facilitador {Perspectivas:}
	Pedagógica, Tecnológica, y Estratégica
\end{itemize}




%\bibliographystyle{plain}
\bibliographystyle{plainurl}
\bibliography{bibliografiaAnexos}

\end{document}
